\section{Introduction}

\subsection{EUV Light Source Design Criteria}

EUVICS is an effort to design and study a compact \EUV{} source based on
\ICS{}. The intended technical program is to establish a traceable relationship
from application requirements through accelerator and laser inputs, interaction
geometry, photon generation and transport, diagnostics, safety interfaces, and
verification evidence. It is not yet a statement of demonstrated wavelength,
flux, bandwidth, brilliance, efficiency, stability, readiness, or experimental
performance.

Application requirements remain approval pending. Candidate requirements
SYS-EUV-001 through SYS-EUV-008 identify wavelength, bandwidth, flux, stability,
polarization, repetition rate, availability, and radiometric traceability, but
their values and verification criteria remain TBD. The 13.5~nm case is a
candidate design target, not an approved requirement. Detector and instrument
calibration must ultimately use traceable, configuration-matched evidence; the
existence of NIST EUV calibration services does not select or qualify an EUVICS
detector~\cite{NISTEUVDetectorCalibration}.

\subsection{Accelerator-Based EUV Light Sources}

An inverse Compton source brings a relativistic electron bunch and an incident
laser pulse into controlled collision. In the electron rest frame, the incident
photon is Doppler blueshifted; it scatters from the electron and is then boosted
again when transformed back to the laboratory frame. This two-Doppler-shift
picture explains the large upshift available from modest incident-photon energy,
while the exact result follows from energy--momentum conservation. High-energy
photons are concentrated about the forward electron direction, but no universal
divergence or collected fraction is assumed: electron phase space,
polarization, collision geometry, observation aperture, and the selected model
must be specified~\cite{SunWu2011ComptonSource}.

\subsection{ICS EUV Light Source Design Concept and Performance Goal}

The nominal roadmap inputs are 3~MeV electron kinetic energy and 800~nm
incident-laser vacuum wavelength. They remain assumed inputs with project-owner
approval pending. The candidate 6.7~nm and 13.5~nm cases are study targets, not
measurements or validated facility performance. A versioned model configuration
must connect any target to electron, laser, geometry, acceptance, optical, and
detector assumptions before the result can be reviewed.

The source chain is: electron production and acceleration; beam transport,
focusing, steering, and diagnostics; laser generation, timing, transport,
polarization control, and focusing; spatiotemporal overlap at the interaction
point; photon scattering; forward EUV collection and spectral selection;
optical transport and calibrated diagnostics; and safe transport or disposal of
the residual electron and laser beams. The interaction point, synchronization,
crossing geometry, and overlap are system interfaces, not secondary corrections.
They determine whether single-particle kinematics translates into useful
accepted and detected output.

\blankline
\blankline
